\section*{第 \ref{ch:g11:plate-motions} 章——板块及其运动}

\begin{solution}{exo:g11:plate-motions:1}
\omterm{def:g11:plate-motions:rigid-plate}{刚性板块}是作为单一整体运动的一块\omterm{def:g11:lithosphere-asthenosphere:lithosphere}{岩石圈}：其上各点之间的距离不变，
其形变集中于沿边界的狭窄带中。地震清楚地显示了这一点：
几乎所有地震都发生在狭窄的条带上，
即板块边界处，
板块内部极少。
\end{solution}

\begin{solution}{exo:g11:plate-motions:2}
\qty{72}{mm} 每年，即每百万年 \qty{72}{km}。30 年后，该台站
移动 $72 \times 30 = \qty{2160}{mm}$，约 \qty{2.2}{m}；百万年后，
\qty{72}{km}。
\end{solution}

\begin{solution}{exo:g11:plate-motions:3}
$v_E = 30/10 = \qty{3}{cm}$ 每年，$v_N = 40/10 = \qty{4}{cm}$ 每年。
速率为 $\sqrt{3^2 + 4^2} = \qty{5}{cm}$ 每年。与正北方向的夹角 $\beta$
满足 $\tan\beta = 3/4$，故 $\beta \approx \ang{37}$：该台站
向东北方向运动，即北偏东 \ang{37}。
\end{solution}

\begin{solution}{exo:g11:plate-motions:4}
\omterm{def:g11:plate-motions:hotspot-track}{热点轨迹}是板块经过热点时，被其带离热点的火山和海山链，
距活火山越远，年龄越老。最年轻的火山基拉韦厄
位于火山链的东南端，在夏威夷岛上；年龄向西西北方向递增。
板块沿其运动方向把火山带离热点：
太平洋板块
向西西北方向
运动。
\end{solution}

\begin{solution}{exo:g11:plate-motions:5}
$519/5.1 \approx \qty{102}{km}$ 每百万年，即 \qty{102}{mm} 每
年，约 \qty{10}{cm} 每年。
\end{solution}

\begin{solution}{exo:g11:plate-motions:6}
$v_E = 63/31 \approx \qty{20.3}{mm}$ 每年，$v_N = 48/31 \approx
\qty{15.5}{mm}$ 每年。速率为 $\sqrt{20.3^2 + 15.5^2} \approx
\qty{25.5}{mm}$ 每年；$\tan\alpha = 15.5/20.3 \approx 0.76$ 得 $\alpha
\approx \ang{37}$：该台站向东偏北 \ang{37} 方向，即向东北方向运动。
茂纳基山的运动速度约快 $72/25.5 \approx 2.8$ 倍。
\end{solution}

\begin{solution}{exo:g11:plate-motions:7}
（1）以毫米每年为单位，
\[
  \vec v_{\mathrm{NA/EU}} = (-10.5 - 7.7,\ 20.6 - 16.3) = (-18.2, +4.3).
\] 其速率为 $\sqrt{18.2^2 + 4.3^2} \approx
\qty{18.7}{mm}$ 每年，$\tan\alpha = 4.3/18.2 \approx 0.24$ 得其方向为
北偏西 \ang{13}。（2）两板块在洋脊处相互远离，速率约为
\qty{1.9}{cm} 每年，即冰岛该洋脊的\omterm{def:g11:drift-to-spreading:spreading-rate}{全扩张速率}。
（3）这与北大西洋的 \qty{2}{cm} 每年
的数值相符。
\end{solution}

\begin{solution}{exo:g11:plate-motions:8}
（1）$v = \pi \times 0.651 \times 6371/180 \approx \qty{72.4}{km}$ 每百万
年，约 \qty{7.2}{cm} 每年。（2）$72.4 \times \sin\ang{30} \approx
\qty{36.2}{km}$ 每百万年，约 \qty{3.6}{cm} 每年。（3）位于自转极处
（若该点位于板块上）。速率随位置而变，
在极点处为零，离极点 \ang{90} 处达到最大值。
因此整个板块并无单一速率，
只有每一点各自的速率（或其表面的平均速率）。
\end{solution}

\begin{solution}{exo:g11:plate-motions:9}
（1）以毫米每年为单位，
\[
  \vec v_{\mathrm{IN/EU}} = (40.6 - 24.0,\ 38.6 + 3.3) = (+16.6, +41.9).
\] 速率为 $\sqrt{16.6^2 + 41.9^2} \approx
\qty{45}{mm}$ 每年；与正北方向的夹角 $\beta$ 满足 $\tan\beta =
16.6/41.9 \approx 0.40$，故 $\beta \approx \ang{22}$：印度向欧亚板块运动，
北偏东 \ang{22}。（2）\qty{45}{km} 每百万年，
故 1000 万年内移动 \qty{450}{km}。
\end{solution}

\begin{solution}{exo:g11:plate-motions:10}
$(2432 - 1058)/(27.7 - 10.3) = 1374/17.4 \approx \qty{79}{km}$ 每百万年，
约 \qty{7.9}{cm} 每年，略小于 \qty{9.1}{cm} 每年。
两火山的年龄各有约五十万年的不确定度；
年龄差 $17.4 \pm 1.0$ 百万年，给出的速率约在 75 至
\qty{84}{km} 每百万年之间。在这两座火山之间，板块的运动似乎
比平均值略慢。
\end{solution}

\begin{solution}{exo:g11:plate-motions:11}
（1）$\vec v_{A/C} = \vec v_{A/B} + \vec v_{B/C}$：向东 \qty{3}{cm} 每年加上
向北 \qty{4}{cm} 每年。两箭头成直角，三角形的
第三边长为 $\sqrt{3^2 + 4^2} = \qty{5}{cm}$ 每年，
方向为北偏东 \ang{37}（$\tan\beta = 3/4$）。（2）$\vec v_{C/A}
= -\vec v_{A/C}$：\qty{5}{cm} 每年，方向相反，
即西偏南 \ang{37}。
\end{solution}

\begin{solution}{exo:g11:plate-motions:12}
（1）$54/0.375 = \qty{144}{km}$ 每百万年，即 \qty{14.4}{cm} 每年。（2）
由 $29/0.375 \approx 77$ 至 $79/0.375 \approx \qty{211}{km}$ 每百万
年：约为 8 至 \qty{21}{cm} 每年。（3）靠近热点处，距离
不大于一座火山的尺寸，或热点向其供给火山的区域的尺寸；
年龄也不大于一座火山生长所需的时间：
二者的相对误差都很大。远离热点处，
\qty{2400}{km} 上 \qty{25}{km} 的误差，
或 2800 万年上五十万年的误差，仅为百分之一或二。
\end{solution}

\begin{solution}{exo:g11:plate-motions:13}
（1）散布的十倍为 \qty{50}{mm}：$50/72 \approx 0.7$ 年，约
八个月。（2）$50/5 = 10$ 年。（3）数百万年间，板块
移动了数十至数百千米，这些距离由海底磁条带记录，
且在海上易于测量；精度达数毫米的定位，
能在数月内观测到运动，出现得晚得多。
\end{solution}

\begin{solution}{exo:g11:plate-motions:14}
（1）以毫米每年为单位，
\[
  \vec v_{\mathrm{SA/AF}} = (-4.9 - 27.1,\ 11.0 - 17.9) = (-32.0, -6.9).
\] 速率为 $\sqrt{32.0^2 + 6.9^2} \approx
\qty{32.7}{mm}$ 每年，由 $\tan\alpha = 6.9/32.0 \approx 0.22$ 得
西偏南 \ang{12}。（2）\omterm{def:g11:drift-to-spreading:spreading-rate}{全扩张速率}约为 \qty{3.3}{cm}
每年，半扩张速率约为 \qty{1.6}{cm} 每年。（3）最近数百万年的半扩张速率
为 \qty{1.6}{cm} 每年，略小于
66 百万年间的平均值 \qty{1.9}{cm} 每年：南大西洋近期的扩张可能比平均情况略慢，
或者差异来自沉积物年龄的不确定度。两种测量结果
一致：
洋中脊的扩张速度很慢。
\end{solution}

\begin{solution}{exo:g11:plate-motions:15}
（1）一条每年上升 \qty{20}{mm} 的直线，
五年后突然下降 \qty{50}{mm}，接着是斜率相同、再次上升的直线。
（2）测站终点位于起点以东 $10 \times 20 - 50 = \qty{150}{mm}$ 处：
$150/10 = \qty{15}{mm}$ 每年，偏小四分之一。（3）两半各得
\qty{20}{mm} 每年。（4）跳跃是地震引起的一次独立的突发运动：
\omterm{def:g10:measuring-the-earth:geodesy}{大地测量学}家根据跳跃前后的位置测量其大小，
将其从序列中去除（或用斜率相同、中间带阶跃的两段拟合），
并取 \qty{20}{mm} 每年，
作为测站的速度。
\end{solution}

\begin{solution}{pb:g11:plate-motions:1}
\textbf{1.} 莫纳克亚火山（37.5 万年）最年轻，须藤海山
（中部）（6470 万年）最古老。年龄随
距基拉韦厄火山的距离增加而增大。

\textbf{2.} $\qty{1}{km} = 10^6~\unit{mm}$，一百万年即 $10^6$
年：每百万年 \qty{1}{km}，即 $10^6~\unit{mm}/10^6$ 年，
\qty{1}{mm} 每年。

\textbf{3.} $519/5.1 \approx \qty{102}{mm}$ 每年。

\textbf{4.} 内克尔岛 $1058/10.3 \approx 102.7$；莱桑岛 $1818/19.9 \approx
91.4$；中途岛 $2432/27.7 \approx 87.8$；科拉汉海山 $3128/38.6 \approx 81.0$；
阿博特海山 $3280/38.7 \approx 84.8$（毫米每年）。

\textbf{5.} 各点 (54; 0.375)、(519; 5.1)、(780; 7.2)、(1058; 10.3)，
(1818; 19.9)、(2432; 27.7)、(3128; 38.6)、(3280; 38.7) 大致落在一条过原点的直线附近，
略有散布（与每百万年 \qty{91}{km} 的直线相比，
考艾岛、尼豪岛和内克尔岛略低于该线，莱桑岛恰在线上，
科拉汉海山和阿博特海山略高于该线）。

\textbf{6.} 每座火山都形成于热点之上，此后随板块一同移动。
设经过的时间为 $t$，板块沿稳定方向移动了 $v\,t$，
因此火山位于距现位于热点上方之点 $d = v\,t$ 处：
即 $t = d/v$。

\textbf{7.} $(102.7 + 91.4 + 87.8 + 81.0 + 84.8)/5 \approx \qty{89.5}{mm}$
每年。

\textbf{8.} 89.5 与每百万年 \qty{91}{km} 相比，二者相差
不到 \qty{2}{\percent}。

\textbf{9.} $519/91 \approx 5.7$ 百万年，而年龄为 \num{5.1}
百万年，二者接近，考艾岛比平均直线
预测的略为年轻。

\textbf{10.} 火山的年龄向西北西方向逐渐增大：
板块沿该方向携带着它们。太平洋板块正朝
西北西方向越过热点移动。

\textbf{11.} $(4860 - 3493)/(64.7 - 42.4) = 1367/22.3 \approx \qty{61}{km}$
每百万年，约合 \qty{6.1}{cm} 每年。

\textbf{12.} $65/91 \approx 0.71$：沿天皇海山链，板块移动的
速度比沿夏威夷海山链约慢 \qty{30}{\percent}。

\textbf{13.} $3668/50 \approx \qty{73}{km}$ 每百万年，约合
\qty{7.3}{cm} 每年，小于每百万年 \qty{91}{km}。过原点的直线
采用 1987 年的年龄，该年龄将夏威夷海山链的终点定在 39 至 43
百万年；较新的、约 50 百万年的转折起点年龄，
则把同样的距离分摊到更长的时间内。若板块在夏威夷海山链早期移动较慢、
后期较快，两种结论便都成立：
速度并非恒定。

\textbf{14.} 太平洋板块运动方向的改变，发生在约
50 百万年前。

\textbf{15.} 链的走向即板块的运动方向，
沿链的年龄间距即其速度。板块可以加速或减速而不转向，
或在不改变速度的情况下转向，二者因此
相互独立。天皇海山链沿线速度变化而方向不变，
说明其速度记录与转折是两件独立的事：转折
并非由热点减速造成。

\textbf{16.} $\sqrt{63.2^2 + 34.5^2} \approx \qty{72.0}{mm}$ 每年，
由 $\tan\alpha = 34.5/63.2 \approx 0.55$ 得方向为西偏北 \ang{29}。

\textbf{17.} 西北西为西偏北 \ang{22.5}，西北
为西偏北 \ang{45}：测站的方向，即西偏北 \ang{29}，位于二者之间，
接近海山链的方向，二者一致。

\textbf{18.} 测站的速度是相对于一个\omterm{def:g11:lithosphere-asthenosphere:lithosphere}{岩石圈}整体不发生转动的参考系测得的，
而海山链给出的是相对于热点的运动：
若两参考系彼此\omterm{def:g11:plate-motions:relative-motion}{相对运动}，则速度不同。
此外，测站测量的是最近 30 年，
海山链对应的是最近 40 百万年，期间速度可能已经改变。
年龄的散布还使海山链的速度带有数毫米
每年的不确定度。

\textbf{19.} $100/91 \approx 1.1$ 百万年。

\textbf{20.} \textbf{太平洋板块在最近 40 百万年间，
以约 \qty{9}{cm} 每年的速度越过夏威夷热点，朝西北西方向移动}
；转折之前它移动得较慢，且
方向不同。
\end{solution}
